% Part: first-order-logic
% Chapter: beyond
% Section: introduction

\documentclass[../../../include/open-logic-section]{subfiles}

\begin{document}

\olfileid{fol}{byd}{int}

\olsection{Overzicht}

De eerste-ordelogica is niet het enige interessante logische systeem:
er bestaan veel uitbreidingen en varianten van. Een logica bestaat
doorgaans uit een formele specificatie van een taal, meestal maar niet
altijd uit een deductiesysteem, en meestal maar niet altijd uit een
bedoelde semantiek. Dit technische gebruik van de term roept echter een
voor de hand liggende vraag op: wat hebben logica's die geen
eerste-ordelogica zijn te maken met het woord ``logica'' in intuïtieve
of filosofische zin? Alle hieronder beschreven systemen zijn bedoeld
om een of andere vorm van redeneren te modelleren; kunnen we aangeven
wat ze logisch maakt?

Een eenvoudig antwoord is niet voorhanden. Het woord ``logica'' wordt
in verschillende betekenissen en contexten gebruikt, en het begrip is,
net als ``waarheid'', vanuit tal van filosofische standpunten
geanalyseerd. Men zou bijvoorbeeld kunnen stellen dat logisch redeneren
tot doel heeft te bepalen welke uitspraken noodzakelijk waar zijn, a
priori waar zijn, waar zijn onafhankelijk van de interpretatie van de
niet-logische termen, waar zijn op grond van hun vorm, of waar zijn
krachtens een taalconventie. Elk van deze opvattingen behoeft veel
nadere uitleg. Zelfs wanneer men zich beperkt tot de logica die in de
wiskunde wordt gebruikt, bestaat weinig overeenstemming over haar
reikwijdte. Zo probeerden Russell en Whitehead in de
\textit{Principia Mathematica} de wiskunde te ontwikkelen op basis van
de logica, in de door Frege begonnen {\em logicistische} traditie. Hun
logische systeem was een vorm van logica met hogere typen, vergelijkbaar
met het hieronder beschreven systeem. Uiteindelijk moesten zij axioma's
invoeren die volgens de meeste maatstaven niet zuiver logisch lijken te
zijn, met name het oneindigheidsaxioma en het axioma van reducibiliteit.
Toch kan men verdedigen dat sommige redeneringen van hogere orde als
logisch moeten gelden. Quine daarentegen, wiens ontologie
``proposities'' niet erkent als legitieme objecten waarover men kan
spreken, betoogt dat tweede- en hogere-ordelogica in feite vermomde
vormen van verzamelingenleer zijn; met andere woorden, systemen waarin
over predicaten wordt gekwantificeerd, zijn niet zuiver logisch.

Voorlopig kunnen we zulke filosofische kwesties het beste voor een
regenachtige dag bewaren en de hieronder beschreven systemen eenvoudig
beschouwen als formele idealiseringen van verschillende soorten
redeneringen, al dan niet logisch.

\end{document}
